循环神经网络，特别是长短期记忆 \citep{hochreiter1997} 和门控循环 \citep{gruEval14} 神经网络，已经被确立为序列建模和转导问题（如语言建模和机器翻译）的最先进方法 \citep{sutskever14, bahdanau2014neural, cho2014learning}。随后许多工作继续推进循环语言模型和编码器-解码器架构的发展 \citep{wu2016google,luong2015effective,jozefowicz2016exploring}。

循环模型通常沿输入和输出序列的符号位置进行计算。将位置与计算时间步对齐，它们生成一个隐态序列 $h_t$，作为前一个隐态 $h_{t-1}$ 和位置 $t$ 的输入的函数。这种固有的顺序性质阻止了在训练示例内的并行化，这在较长的序列长度下变得至关重要，因为内存约束限制了示例间的批处理。
%\marginpar{not sure if the memory constraints are understandable here}
最近的工作通过分解技巧 \citep{Kuchaiev2017Factorization} 和条件计算 \citep{shazeer2017outrageously} 在计算效率上取得了显著改进，同时也改进了后者的模型性能。然而，顺序计算的根本约束仍然存在。

%\marginpar{@all: there is work on analyzing what attention really does in seq2seq models, couldn't find it right away} 

注意力机制已成为各种任务中令人信服的序列建模和转导模型的重要组成部分，允许在不考虑其在输入或输出序列中的距离的情况下对依赖关系进行建模 \citep{bahdanau2014neural, structuredAttentionNetworks}。然而，除少数情况外 \citep{decomposableAttnModel}，这样的注意力机制都是与循环网络结合使用的。

%\marginpar{not sure if "cross-positional communication" is understandable without explanation}
%\marginpar{insert exact training times and stats for the model that reaches sota earliest, maybe even a single GPU model?}

在本工作中，我们提出了 Transformer，一种避免循环而完全依赖注意力机制来捕捉输入和输出之间全局依赖的模型架构。Transformer 允许显著更多的并行化，并且在仅用八块 P100 GPU 训练十二小时后就能达到翻译质量的新最优水平。 
%\marginpar{you removed the constant number of repetitions part. I wrote it because I wanted to make it clear that the model does not only perform attention once, while it's also not recurrent. I thought that might be important to get across early.}

% Just a standard paragraph with citations, rewrite.
%After the seminal papers of \citep{sutskever14}, \citep{bahdanau2014neural}, and \citep{cho2014learning}, recurrent models have become the dominant solution for both sequence modeling and sequence-to-sequence transduction. Many efforts such as \citep{wu2016google,luong2015effective,jozefowicz2016exploring} have pushed the boundaries of machine translation and language modeling with recurrent sequence models. Recent effort \citep{shazeer2017outrageously} has combined the power of conditional computation with sequence models to train very large models for machine translation, pushing SOTA at lower computational cost. Recurrent models compute a vector of hidden states $h_t$, for each time step $t$ of computation. $h_t$ is a function of both the input at time $t$ and the previous hidden state $h_t$. This dependence on the previous hidden state encumbers recurrnet models to process multiple inputs at once, and their time complexity is a linear function of the length of the input and output, both during training and inference. [What I want to say here is that although this is fine during decoding, at training time, we are given both input and output and this linear nature does not allow the RNN to process all inputs and outputs simultaneously and haven't been used on datasets that are the of the scale of the web. What's the largest dataset we have ? . Talk about Nividia and possibly other's effors to speed up things, and possibly other efforts that alleviate this, but are still limited by it's comptuational nature]. Rest of the intro: What if you could construct the state based on the actual inputs and outputs, then you could construct them all at once. This has been the foundation of many promising recent efforts, bytenet,facenet (Also talk about quasi rnn here). Now we talk about attention!! Along with cell architectures such as long short-term meory (LSTM) \citep{hochreiter1997}, and gated recurrent units (GRUs) \citep{cho2014learning}, attention has emerged as an essential ingredient in successful sequence models, in particular for machine translation. In recent years, many, if not all, state-of-the-art (SOTA) results in machine translation have been achieved with attention-based sequence models \citep{wu2016google,luong2015effective,jozefowicz2016exploring}. Talk about the neon work on how it played with attention to do self attention! Then talk about what we do.